\section{Background and Motivation}
\label{sec:background}

We examine why existing host-only and off-path receive architectures fall short
at 100--800\,Gbps, highlighting two hardware realities---the pre-posted DMA
ring lower bound and coarse hardware backpressure---that shape \sys{}.

\subsection{DDIO and the Pre-Posted Ring Lower Bound}
\label{sec:bg-ddio}

On modern servers, Intel DDIO~\cite{intelddio,farshin2020ddio} places inbound
NIC DMA writes directly into a small subset of LLC ways (typically 2 ways, or
$\sim$8\,MiB out of 60\,MiB on our Xeon Sapphire Rapids socket) when target
lines miss, or updates them in place when already cached. As long as active
receive buffers remain inside this budget, packets are consumed at LLC speed
with minimal DRAM traffic. Once active buffers outgrow DDIO, new DMA writes
evict unread packets to DRAM (``leaky DMA'')~\cite{tootoonchian2018resq,farshin2020ddio},
amplifying host memory traffic and inflating tail latency for co-located
tenants.

\begin{table*}[t]
\centering
\footnotesize
\setlength{\tabcolsep}{4.0pt}
\renewcommand{\arraystretch}{1.08}
\begin{tabularx}{\textwidth}{l c c c c c Y}
\toprule
\textbf{Architecture} &
\textbf{Burst Buffer} &
\textbf{Lock-Free?} &
\textbf{Slow-Core} &
\textbf{Host API / DMA} &
\textbf{Flow State} &
\textbf{Structural Limitation Addressed by \sys{}} \\
\midrule
\texttt{privRing} & Host RAM & Yes & Per-queue & Standard & None &
$N \times D \times \text{MTU}$ blows DDIO; leaky DMA degrades co-located tenants. \\
\textbf{ShRing}~\cite{pismenny2023shring} & Host RAM & No & HoL / fallback & Standard & None &
Lock contention; stalled core blocks peers or forces fallback to \texttt{privRing}. \\
\textbf{rxBisect}~\cite{pismenny2025rxbisect} & Host RAM ($B_x$) & Yes & Bounded ($B_x$) & Custom ASIC & None &
Host-memory conservation wall: large bursts/stalls overflow $B_x$ or bloat host LLC. \\
\textbf{\ceio{}}~\cite{liu2025ceio} & Off-path NIC & Yes & Slow path & Modified (\texttt{async}) & $O(F)$ RMT &
Dual-path tax: RMT lag, phase pause, 2P1C SW merge, fast-pool LLC read contention, churn drop. \\
\midrule
\textbf{\sys{}} & \textbf{On-path NIC} & \textbf{Yes} & \textbf{Full (NIC)} & \textbf{Standard DMA} & \textbf{$O(1)$/Q} &
Single unified path; normal DPDK DMA + compact pool + unfenced 64\,B MMIO credits. \\
\bottomrule
\end{tabularx}
\caption{Evolution of cache-efficient NIC receive-ring architectures from private rings to \sys{}.}
\label{tab:design-space}
\end{table*}

Replacing a host FIFO buffer pool with an ideal LIFO stack (\texttt{LargeLifo})
while keeping large hardware rings cannot fix leaky DMA. Because a standard DMA
NIC consumes pre-posted descriptors strictly in FIFO order, any cache-hot
\texttt{mbuf} returned to a software LIFO stack is re-posted to the \emph{tail}
of a $D$-entry hardware ring and is not overwritten until $N_q \cdot D$
intervening packets rotate through all $N_q$ rings. The effective DMA working
set therefore obeys a strict lower bound:
\begin{equation}
  W_{\text{eff}} = \max\!\bigl(N_q \cdot D,\; N_{\text{burst}}\bigr) \cdot \text{sizeof}(\text{mbuf}).
  \label{eq:fifo-lower-bound}
\end{equation}
With four 8192-descriptor rings and 2304\,B \texttt{mbuf}s (1500\,B frames),
pre-posted rings alone span $32{,}768 \times 2304\,\text{B} = 72.0$\,MiB---exceeding
the entire 60\,MiB LLC so even a 375\,KiB burst collapses DDIO to 2.0\%
(\S\ref{sec:eval-e2}). Achieving DDIO residency therefore requires shrinking
\emph{both} the hardware descriptor ring $D$ and the backing buffer pool.

\subsection{Limits of Host-Only and Dual-Path Designs}
\label{sec:bg-evolution}

Table~\ref{tab:design-space} contrasts four architectural generations.
\textbf{ShRing}~\cite{pismenny2023shring} shares one compact ring across cores
without NIC changes, but adds cross-core synchronization and allows a single
stalled core to pin the ring head, starving healthy cores until ShRing falls
back to \texttt{privRing} and forfeits cache isolation.
\textbf{rxBisect}~\cite{pismenny2025rxbisect} removes locks by pairing a small
lock-free shared allocation ring ($A_x$) with large per-core delivery rings
($B_x$), yet still buffers all backlog inside host RAM: under 5--10\,ms stalls
at 100\,Gbps (10K--45K packets), bounded $B_x$ overflows and drops packets,
while enlarging $B_x$ evicts co-located workloads from LLC---in addition to
requiring custom NIC ASIC modifications.

\textbf{\ceio{}}~\cite{liu2025ceio} moves burst buffering to off-path SmartNIC
DRAM (Figure~\ref{fig:arch}(a)), but bypassing sidecar ARM cores on the fast
path forces reactive dual-path switching. As we quantify head-to-head in
\S\ref{sec:eval-e3}, this introduces four persistent penalties: (1)~polling
separate fast/slow rings and merging via a 2P1C software ring adds
$+74.5$\,cycles/packet even at 0\% slow-path traffic; (2)~per-flow credit
accounting and RMT table limits degrade throughput under thousands of churning
flows; (3)~oversized fast-path pools (2048 slots/queue) needed to absorb
switching overshoot still contend with co-located tenants for shared LLC read
capacity despite $>$99.5\% DDIO; and (4)~applications must adopt a new
\texttt{async\_recv()} pull interface.

\subsection{Why Native NIC Hardware Backpressure Falls Short}
\label{sec:bg-backpressure}

Shrinking a standard DMA ring and relying solely on native NIC hardware
backpressure is also insufficient. First, hardware descriptor watermarks are
coarse (on CN106 \texttt{octeon\_ep}, $D \ge 256$, $W_{\text{mark}} \ge 64$,
refill stride 64) and cannot enforce fine packet caps (e.g., $R_{\text{eff}}=64$)
or byte-level LLC budgets (64--96\,KiB). Second, hardware backpressure risks
multi-queue HoL stalling in shared NIC egress pipelines. Third, on CN106,
enabling SDP-to-NIX channel backpressure requires setting
\texttt{SDPX\_OUT\_BP\_ENX\_W1S}, which the Linux
kernel driver sets but the default user-space endpoint agent
(\texttt{octep\_cp\_agent}) leaves cleared; consequently, default 256--512-slot
rings silently drop 410K--1.34M packets under host pauses without engaging
SmartNIC memory (\S\ref{sec:eval-e1}). While silent non-accounting
(\texttt{imissed=0}) is specific to CN106's user-space stack, passive small
rings overflow universally, making explicit software-visible credits essential
for portable, exact, per-queue isolation.
